\documentclass[12pt,a4paper]{article}

% --- Paquetes Esenciales ---
\usepackage[utf8]{inputenc}
\usepackage[T1]{fontenc}
\usepackage{amsmath}
\usepackage{amssymb}
\usepackage{graphicx}
\usepackage[spanish]{babel}
\usepackage[a4paper, margin=2.5cm]{geometry}
\usepackage{comment}

% --- Paquetes para Mejoras Visuales ---
\usepackage{hyperref} % Para enlaces clickeables
\hypersetup{
    colorlinks=true,
    linkcolor=blue,
    filecolor=magenta,      
    urlcolor=cyan,
}
\usepackage{xcolor} % Para definir colores
\usepackage{listings} % Para bloques de código
\usepackage{framed} % Para enmarcar las preguntas

% --- Definición de Colores y Estilos ---
\definecolor{codegreen}{rgb}{0,0.6,0}
\definecolor{codegray}{rgb}{0.5,0.5,0.5}
\definecolor{codepurple}{rgb}{0.58,0,0.82}
\definecolor{backcolour}{rgb}{0.95,0.95,0.92}

\lstdefinestyle{mystyle}{
    backgroundcolor=\color{backcolour},   
    commentstyle=\color{codegreen},
    keywordstyle=\color{magenta},
    numberstyle=\tiny\color{codegray},
    stringstyle=\color{codepurple},
    basicstyle=\ttfamily\footnotesize,
    breakatwhitespace=false,         
    breaklines=true,                 
    captionpos=b,                    
    keepspaces=true,                 
    numbers=left,                    
    numbersep=5pt,                  
    showspaces=false,                
    showstringspaces=false,
    showtabs=false,                  
    tabsize=2
}
\lstset{style=mystyle}

\newenvironment{pregunta}
  {\begin{leftbar}\begin{quote}\itshape}
  {\end{quote}\end{leftbar}}

\title{Resumen de Conversación sobre Diseño de Circuitos Integrados y Herramientas EDA}
\author{Gemini Code Assist y Prof. CERO}
\date{\today}

\begin{document}

\maketitle
\tableofcontents
\newpage

\section{¿Qué es The OpenROAD Project?}

\begin{pregunta}
Saludos Gemini, que sabes del programa 'Open Road' que sirve para el diseño de circuitos integrados?
\end{pregunta}

¡Hola! Claro que sí. Conozco bien el proyecto al que te refieres. Su nombre correcto es \textbf{The OpenROAD Project}, y es una iniciativa muy importante en el mundo del diseño de circuitos integrados.

Es un conjunto de herramientas de software de \textbf{código abierto (open-source)} para la automatización del diseño electrónico (EDA, por sus siglas en inglés). Su objetivo principal es proporcionar un \textbf{flujo de diseño digital completamente autónomo}, desde el RTL (Register-Transfer Level) hasta el GDSII.

En términos más sencillos, OpenROAD busca automatizar todo el proceso de diseño físico de un chip digital, tomando como entrada una descripción del circuito en un lenguaje como Verilog y produciendo como salida el archivo final (GDSII) que se envía a la fábrica (la \textit{foundry}) para su manufactura.

\subsection{¿Cuál es su principal objetivo?}

El lema del proyecto es \textbf{'No-Human-in-the-Loop'}, lo que significa que el flujo de diseño está pensado para funcionar de manera completamente automática, sin necesidad de intervención manual por parte de un ingeniero. La meta es poder completar el diseño de un chip en 24 horas.

Esto es revolucionario porque los flujos de diseño tradicionales, que utilizan herramientas comerciales muy costosas (de empresas como Synopsys, Cadence o Siemens EDA), a menudo requieren semanas o meses de trabajo y la intervención de múltiples ingenieros expertos para ajustar y optimizar cada etapa del proceso.

\subsection{Etapas del Flujo de Diseño que Cubre OpenROAD}

OpenROAD integra una serie de herramientas, cada una especializada en una fase del diseño físico del chip:

\begin{enumerate}
    \item \textbf{Síntesis Lógica:} Transforma el código RTL (Verilog) en una lista de compuertas lógicas estándar (netlist). A menudo utiliza herramientas externas como \texttt{Yosys}.
    \item \textbf{Floorplanning:} Define la distribución general del chip, la ubicación de los bloques principales (macros), los pines de entrada/salida y la red de alimentación.
    \item \textbf{Placement (Ubicación):} Coloca las células estándar (compuertas lógicas) en las filas definidas durante el floorplanning.
    \item \textbf{Clock Tree Synthesis (CTS):} Construye la red que distribuye la señal de reloj a todos los flip-flops del diseño, asegurando que llegue a todos los puntos al mismo tiempo.
    \item \textbf{Routing (Enrutado):} Conecta todas las compuertas y bloques entre sí mediante pistas metálicas en las diferentes capas del chip.
    \item \textbf{Análisis y Verificación:} Realiza análisis de tiempo estático (STA) para verificar que el chip cumple con las especificaciones de velocidad, y comprobaciones de reglas de diseño (DRC) para asegurar que el layout es fabricable.
\end{enumerate}

\subsection{¿Por qué es tan importante?}

\begin{itemize}
    \item \textbf{Democratización del Diseño de Chips:} Al ser de código abierto y gratuito, permite que universidades, startups, investigadores y entusiastas puedan diseñar circuitos integrados complejos sin tener que pagar licencias de software que cuestan cientos de miles o millones de dólares.
    \item \textbf{Innovación y Transparencia:} Al ser open-source, permite a la comunidad académica y a la industria investigar, modificar y mejorar los algoritmos de diseño. Esto fomenta la innovación de una manera que no es posible con las herramientas comerciales 'de caja negra'.
    \item \textbf{Apoyo Institucional:} Es un proyecto financiado por DARPA (la Agencia de Proyectos de Investigación Avanzados de Defensa de EE. UU.), lo que le da un gran respaldo y credibilidad.
\end{itemize}

\newpage

\section{Primeros Pasos con OpenROAD}

\begin{pregunta}
¿Cómo puedo empezar a usar The OpenROAD Project para un diseño simple?
\end{pregunta}

La forma más sencilla y recomendada para empezar es usando \textbf{Docker}, ya que te abstrae de la compleja instalación de dependencias.

\subsection{Paso 1: Prerrequisitos}

\begin{itemize}
    \item Docker instalado y funcionando.
    \item Un editor de texto.
    \item Conocimientos básicos de Verilog.
\end{itemize}

\subsection{Paso 2: Preparar el Entorno y el Diseño}

Vamos a crear un diseño muy simple: un inversor.

\subsubsection{Archivo de diseño Verilog (\texttt{inversor.v})}
\begin{lstlisting}[language=Verilog]
module inversor (
    input  a,
    output y
);
    assign y = ~a;
endmodule
\end{lstlisting}

\subsubsection{Archivo de restricciones de diseño (\texttt{constraints.sdc})}
\begin{lstlisting}[language=TCL]
# 1. Definir la frecuencia del reloj (ej. 500 MHz -> 2.0 ns de periodo)
create_clock -name clk -period 2.0 [get_ports clk]

# 2. Especificar incertidumbre del reloj (para setup y hold)
set_clock_uncertainty 0.1 [get_clocks clk]

# 3. Definir los retardos de entrada y salida
set_input_delay 0.5 -clock clk [all_inputs]
set_output_delay 0.5 -clock clk [all_outputs]
\end{lstlisting}

\subsubsection{Archivo de configuración (\texttt{config.tcl})}
\begin{lstlisting}[language=TCL]
# 1. Plataforma y tecnologia
set ::env(PLATFORM) 'nangate45'

# 2. Nombre del diseno (debe coincidir con el 'module' en Verilog)
set ::env(DESIGN_NAME) 'inversor'

# 3. Archivos de entrada Verilog
set ::env(VERILOG_FILES) '$::env(DESIGN_DIR)/inversor.v'

# 4. Archivo de restricciones (SDC)
set ::env(SDC_FILE) '$::env(DESIGN_DIR)/constraints.sdc'

# 5. Parametros del Floorplan (tamano del chip)
set ::env(FP_CORE_UTIL) 50
set ::env(FP_ASPECT_RATIO) 1.0

# 6. Parametros del reloj
set ::env(CLOCK_PERIOD) '2.0' ;# ns
set ::env(CLOCK_PORT) 'clk'
\end{lstlisting}

\subsection{Paso 3: Ejecutar el Flujo de OpenROAD}

Primero, descarga la imagen de Docker:
\begin{lstlisting}[language=bash]
docker pull openroad/openroad
\end{lstlisting}

Luego, ejecuta el flujo:
\begin{lstlisting}[language=bash]
docker run -it --rm \
    -v $(pwd):/openroad/designs \
    -e DESIGN_DIR=/openroad/designs \
    -u $(id -u):$(id -g) \
    openroad/openroad \
    bash -c 'source /openroad/etc/openroad_init && flow -design /openroad/designs'
\end{lstlisting}

\subsection{Paso 4: Revisar los Resultados}

Al finalizar, tu directorio contendrá una nueva carpeta llamada \texttt{results}. Los archivos más importantes son:
\begin{itemize}
    \item \texttt{results/nangate45/inversor/6\_final.gds}: ¡Este es el archivo final! Es el layout que enviarías a fabricar.
    \item \texttt{results/nangate45/inversor/6\_final.def}: El diseño final en formato DEF, que puedes visualizar.
    \item \texttt{logs/}: Contiene los logs detallados de cada paso.
    \item \texttt{reports/}: Contiene informes de tiempo, área, potencia, etc.
\end{itemize}

\newpage
\section{Diferencias entre LEF, DEF y GDS}

\begin{pregunta}
Explícame las diferencias entre los archivos LEF, DEF y GDS en el contexto de OpenROAD.
\end{pregunta}

Entender el propósito de estos tres formatos de archivo es fundamental. Son como los planos en diferentes etapas de la construcción de un edificio.

\subsection{LEF (Library Exchange Format) - La 'Librería de Componentes'}
\begin{itemize}
    \item \textbf{Propósito:} Es un archivo de \textbf{abstracción física}. Describe las propiedades físicas de los 'bloques de construcción' que usarás, pero \textbf{NO} dice dónde se colocan ni cómo se conectan.
    \item \textbf{Analogía:} La librería de componentes para diseñar un PCB. El LEF es como el archivo que define el \textit{footprint} de un componente.
    \item \textbf{Rol en OpenROAD:} Es un \textbf{archivo de entrada} fundamental.
\end{itemize}

\subsection{DEF (Design Exchange Format) - El 'Plano de Colocación y Cableado'}
\begin{itemize}
    \item \textbf{Propósito:} Describe la \textbf{implementación física de TU diseño específico}. Utiliza los componentes definidos en el LEF y especifica exactamente dónde se coloca cada uno y cómo se conectan.
    \item \textbf{Analogía:} El archivo del proyecto de tu PCB ya terminada, mostrando la ubicación de cada componente y las pistas que los unen.
    \item \textbf{Rol en OpenROAD:} Es un formato de \textbf{salida intermedia y final}.
\end{itemize}

\subsection{GDSII (Graphic Data System) - El 'Archivo para Fabricación'}
\begin{itemize}
    \item \textbf{Propósito:} Es el formato final, el 'plano maestro' que se envía a la fábrica. Es una representación puramente geométrica del diseño.
    \item \textbf{Analogía:} Un conjunto de archivos Gerber para un PCB. No contiene información lógica, solo polígonos, rectángulos y trazados en diferentes capas.
    \item \textbf{Rol en OpenROAD:} Es la \textbf{salida final y definitiva} del flujo.
\end{itemize}

\subsection{Resumen Comparativo}
\begin{center}
\begin{tabular}{|l|p{4cm}|p{4cm}|p{4cm}|}
\hline
\textbf{Característica} & \textbf{LEF} & \textbf{DEF} & \textbf{GDSII} \\
\hline
\textbf{Propósito} & Librería de componentes físicos & Plano de implementación del diseño & Plano final para fabricación \\
\hline
\textbf{Contenido} & Vistas abstractas de celdas y reglas & Instancias, ubicaciones y rutas & Geometrías (polígonos) por capa \\
\hline
\textbf{Nivel de Abstracción} & Alto (Define los 'qué') & Medio (Define los 'dónde' y 'cómo') & Bajo (Define las 'formas') \\
\hline
\textbf{Rol en el Flujo} & \textbf{Entrada} & Salida intermedia y final & \textbf{Salida final} \\
\hline
\textbf{Legibilidad} & Texto, legible & Texto, legible & Binario, no legible directamente \\
\hline
\end{tabular}
\end{center}

\newpage
\section{Herramientas de Visualización}

\begin{pregunta}
¿Qué herramientas de código abierto puedo usar para visualizar archivos LEF, DEF y GDSII?
\end{pregunta}

\begin{description}
    \item[KLayout (La opción más recomendada):] Es considerado el 'estándar de oro' de los visualizadores de layouts de código abierto. Es extremadamente potente, rápido y versátil. Soporta GDSII, LEF/DEF y muchos otros. Ideal para inspección final y visualización general.
    
    \item[Magic:] Una de las herramientas EDA de código abierto más antiguas y veneradas. Es un editor de layout VLSI completo, no solo un visor. Ideal para edición de bajo nivel y desarrollo de celdas personalizadas.
    
    \item[OpenROAD GUI:] El propio The OpenROAD Project incluye una interfaz gráfica para visualizar el estado del diseño en cada etapa del flujo. Es la mejor herramienta para depurar y entender por qué la herramienta tomó ciertas decisiones en una etapa específica.
\end{description}

\section{Scripts de Análisis con KLayout}

\begin{pregunta}
Muéstrame un ejemplo de script de KLayout en Python para analizar un archivo GDSII.
\end{pregunta}

Este script calcula el área total ocupada por el metal en una capa específica.
\begin{lstlisting}[language=Python]
# Importa el modulo 'pya', que es la API de Python para KLayout
import pya
import os

# --- 1. Configuracion ---
GDS_FILE = 'mi_diseno_ejemplo.gds'
LAYER_TO_ANALYZE = 1
DATATYPE_TO_ANALYZE = 0

# --- 2. Funcion para crear un GDS de prueba (si no existe) ---
def crear_gds_de_prueba(filepath, layer, datatype):
    # ... (codigo para crear un GDS de ejemplo)
    pass

# --- 3. Funcion principal de analisis ---
def analizar_area_de_capa(gds_file, layer_num, datatype_num):
    try:
        layout = pya.Layout()
        layout.read(gds_file)
        top_cell = layout.top_cell()
        layer_info = pya.LayerInfo(layer_num, datatype_num)
        layer_index = layout.layer(layer_info)
        total_area = 0.0
        for shape in top_cell.shapes(layer_index).each():
            total_area += shape.area()
        dbu = layout.dbu
        area_in_um2 = total_area * dbu * dbu
        print(f'Area total calculada: {area_in_um2:.4f} um^2')
    except Exception as e:
        print(f'Ocurrio un error: {e}')

# --- 4. Ejecucion del script ---
if __name__ == '__main__':
    # crear_gds_de_prueba(GDS_FILE, LAYER_TO_ANALYZE, DATATYPE_TO_ANALYZE)
    # analizar_area_de_capa(GDS_FILE, LAYER_TO_ANALYZE, DATATYPE_TO_ANALYZE)
    print('Ejemplo de script de KLayout para analisis de area.')
\end{lstlisting}

\begin{pregunta}
Explícame cómo realizar una operación booleana (AND, OR, XOR) entre dos capas usando un script de KLayout.
\end{pregunta}

En KLayout, se utilizan objetos \texttt{pya.Region} y sus operadores sobrecargados (\texttt{\&} para AND, \texttt{|} para OR, \texttt{\^} para XOR).
\begin{lstlisting}[language=Python]
import pya

def ejecutar_operaciones_booleanas(input_file, output_file):
    layout = pya.Layout()
    layout.read(input_file)
    top_cell = layout.top_cell()
    
    # Define capas de entrada y salida
    layer1_info = pya.LayerInfo(1, 0)
    layer2_info = pya.LayerInfo(2, 0)
    layer_and_info = pya.LayerInfo(10, 0)
    
    # Crea objetos 'Region' para las capas de entrada
    region1 = pya.Region(top_cell.shapes(layout.layer(layer1_info)))
    region2 = pya.Region(top_cell.shapes(layout.layer(layer2_info)))

    # Realiza la operacion booleana AND
    region_and = region1 & region2

    # Inserta la region resultante en la capa de salida
    top_cell.shapes(layout.layer(layer_and_info)).insert(region_and)
    
    layout.write(output_file)
\end{lstlisting}

\newpage
\section{Bibliografía y Aprendizaje}

\begin{pregunta}
Que bibliografía me recomiendas para aprender a diseña circuitos integrados con esta herramienta 'Open Road'?
\end{pregunta}

Dado que es una herramienta moderna, la 'bibliografía' tradicional es escasa. Las mejores fuentes son recursos en línea.

\subsection{Conocimientos Previos Fundamentales}
\begin{itemize}
    \item \textbf{Diseño Lógico Digital:} 'Digital Design' por M. Morris Mano.
    \item \textbf{Lenguajes de Descripción de Hardware (HDL):} 'Verilog HDL: A Guide to Digital Design and Synthesis' por Samir Palnitkar.
    \item \textbf{Diseño Físico de Circuitos VLSI (Conceptos):} 'Digital Integrated Circuits: A Design Perspective' por Jan M. Rabaey.
\end{itemize}

\subsection{Recursos Directos para Aprender OpenROAD}
\begin{description}
    \item[VLSI System Design (VSD) Courses (MUY RECOMENDADO):] Ofrecen talleres y cursos en línea, muchos gratuitos, que guían paso a paso en el uso de herramientas EDA de código abierto.
    \item[The OpenROAD-flow-scripts Tutorial:] El tutorial oficial en el repositorio de GitHub. Es la fuente de verdad.
    \item[Documentación Oficial de OpenROAD (Read the Docs):] Ayuda a entender qué hace cada comando y cada herramienta individualmente.
    \item[Comunidad y Soporte (Slack, GitHub):] La mejor manera de resolver dudas específicas es preguntar a la comunidad.
\end{description}

\newpage
\section{De lo Analógico a lo Digital}

\begin{pregunta}
Estas en capacidad de ayudarme a diseña un circuito integrado?
\end{pregunta}

Sí, por supuesto. Aunque no puedo reemplazar el trabajo de un ingeniero, soy una herramienta excelente para asistirte en muchas de las tareas clave del proceso, como:
\begin{itemize}
    \item Análisis y diseño de circuitos.
    \item Escritura y depuración de código Verilog.
    \item Uso del flujo de diseño con OpenROAD.
    \item Generación y análisis de archivos de layout.
    \item Conceptualización y planificación de la arquitectura.
\end{itemize}

\begin{pregunta}
¿Cómo puedo convertir un circuito analógico simple en una descripción que pueda usar en un flujo de diseño digital?
\end{pregunta}

Un flujo de diseño digital como el de OpenROAD no puede procesar directamente componentes analógicos. La solución es crear un \textbf{modelo digital} del comportamiento del circuito analógico o \textbf{interactuar} con él a través de una interfaz.

\subsection{Enfoque 1: Modelar el Comportamiento (El Enfoque del 'Datapath')}

Consiste en crear un circuito digital que \textbf{calcula el resultado} de las ecuaciones del circuito analógico. Por ejemplo, para la fórmula $V_A = (E1 \cdot R2 + E2 \cdot R1) / (R1 + R2)$:
\begin{enumerate}
    \item \textbf{Digitalizar las Entradas:} Las tensiones de entrada $E1$ y $E2$ deben ser convertidas a números digitales usando un Convertidor Analógico-Digital (ADC).
    \item \textbf{Crear un Datapath en Verilog:} Se escribe código Verilog que implemente la fórmula matemática usando los valores digitales.
    \item \textbf{Ejecutar el Flujo Digital:} Este archivo Verilog se usa como entrada para OpenROAD, que construirá un circuito de compuertas (sumadores, multiplicadores, etc.).
\end{enumerate}

\subsection{Enfoque 2: Abstraer a un Nivel Lógico (El Enfoque del 'Comparador')}

Se usa cuando solo importa si una tensión está por encima o por debajo de un cierto umbral.
\begin{enumerate}
    \item \textbf{Usar un Comparador:} Es un circuito analógico simple cuya salida es digital ('1' o '0').
    \item \textbf{Conectar la Salida al Flujo Digital:} La salida del comparador se trata como una entrada digital simple en el diseño Verilog.
\end{enumerate}

\newpage
\section{Fundamentos del Diseño Digital}

\begin{pregunta}
Entonces, para el diseño de un circuito integrado se comienza haciendo un programa en el lenguaje verilog?
\end{pregunta}

¡Exacto! Pero con una distinción importante: no es un 'programa', es una \textbf{'descripción'}.
\begin{itemize}
    \item Un \textbf{programa} define una \textbf{secuencia de instrucciones} que se ejecutan una tras otra.
    \item Una \textbf{descripción en Verilog} define una \textbf{estructura de hardware}: un conjunto de componentes y conexiones que existen \textbf{simultáneamente} y operan en paralelo.
\end{itemize}

El flujo de diseño digital comienza así:
\begin{enumerate}
    \item \textbf{Especificación y Arquitectura:} Diseño 'en papel'.
    \item \textbf{Diseño RTL (Register-Transfer Level):} Se escribe el código Verilog.
    \item \textbf{Simulación y Verificación Funcional:} Se prueba el código Verilog con un banco de pruebas (testbench).
    \item \textbf{Síntesis y Diseño Físico (El Flujo de OpenROAD):} El Verilog verificado se entrega al flujo de automatización para crear el layout físico.
\end{enumerate}

\begin{pregunta}
Que es un 'datapath digital'?
\end{pregunta}

Un \textbf{datapath digital} (o ruta de datos) es la parte de un procesador o sistema digital que contiene los bloques funcionales que realizan las operaciones de procesamiento de datos. Es el 'músculo' del sistema.

Está compuesto por:
\begin{itemize}
    \item \textbf{Elementos de Procesamiento:} Realizan operaciones sobre los datos.
    \begin{itemize}
        \item \textbf{Unidad Aritmético-Lógica (ALU):} Realiza operaciones como suma, resta, AND, OR, etc.
        \item \textbf{Multiplicadores, Divisores, Shifters.}
    \end{itemize}
    \item \textbf{Elementos de Almacenamiento:} Guardan los datos temporalmente.
    \begin{itemize}
        \item \textbf{Registros (Registers):} Guardan resultados intermedios.
        \item \textbf{Buses:} Los 'cables' que conectan todo.
        \item \textbf{Multiplexores (MUX):} Actúan como 'interruptores' para seleccionar qué dato, de entre varias fuentes, se envía a un bloque funcional.
    \end{itemize}
\end{itemize}
El datapath es controlado por la \textbf{Unidad de Control}, que es el 'cerebro' que le dice al datapath qué operación ejecutar.

\begin{pregunta}
¿Cuál es la diferencia entre un datapath y una ALU?
\end{pregunta}

La diferencia es de escala y función, como la diferencia entre un motor y un coche completo. En resumen: \textbf{La ALU es un componente \textit{dentro} del datapath.}

\begin{description}
    \item[ALU (Unidad Aritmético-Lógica):] Es el 'calculador' del sistema. Su única tarea es tomar dos operandos de entrada y una señal de control para producir un resultado. Es como una calculadora de bolsillo.
    \item[Datapath (Ruta de Datos):] Es la infraestructura completa que permite el flujo y la manipulación de datos. Contiene a la ALU, además de Registros, Multiplexores y Buses. Es como el escritorio de trabajo completo de un contable, que incluye la calculadora, las hojas de papel y sus manos para mover la información.
\end{description}

\newpage
\section{El Ecosistema Abierto: RISC-V y OpenROAD}

\begin{pregunta}
Ok, siguiendo en el contexto de diseño de circuitos integrados con OpenROAD project, hablame de RISC-V.
\end{pregunta}

Hablar de RISC-V es el siguiente paso lógico y natural. RISC-V y The OpenROAD Project son las dos piezas clave que están impulsando una verdadera revolución en el hardware de código abierto.

Piénsalo de esta manera:
\begin{itemize}
    \item \textbf{The OpenROAD Project} te da las \textit{herramientas} para construir un chip (la fábrica de software).
    \item \textbf{RISC-V} te da el \textit{plano} para construir el cerebro del chip (el procesador).
\end{itemize}

Juntos, eliminan las dos barreras más grandes y costosas para el diseño de circuitos integrados personalizados.

\subsection{¿Qué es RISC-V?}

\textbf{RISC-V} (pronunciado 'risk-five') es una \textbf{Arquitectura de Conjunto de Instrucciones (ISA)} de código abierto.

\begin{description}
    \item[ISA (Instruction Set Architecture):] Es la especificación fundamental de un procesador. Define las instrucciones que el procesador puede ejecutar (como sumar, cargar datos de memoria, etc.), los registros que tiene y cómo maneja la memoria. Es el contrato entre el hardware y el software. Otros ejemplos de ISAs son x86 (de Intel/AMD) y ARM.
    
    \item[Código Abierto (Open Standard):] A diferencia de x86 y ARM, que son propietarias y requieren licencias muy costosas, la especificación de RISC-V es \textbf{completamente abierta y de uso gratuito}. Cualquiera puede descargar la especificación y diseñar su propio procesador RISC-V sin pagar regalías ni pedir permiso.
\end{description}

\textbf{Punto clave:} RISC-V no es un procesador. Es un \textbf{estándar} o una \textbf{receta} para construir procesadores. Hay muchas implementaciones (llamadas 'cores' o núcleos) de procesadores que siguen la receta de RISC-V, algunas comerciales y muchas de código abierto.

\subsection{La Sinergia Perfecta: RISC-V + OpenROAD}

Aquí es donde la magia ocurre. La combinación de un estándar de procesador abierto (RISC-V) y un flujo de herramientas de diseño abierto (OpenROAD) crea un ecosistema increíblemente poderoso:

\begin{enumerate}
    \item \textbf{Democratización Total:} Antes, para diseñar un System-on-Chip (SoC), necesitabas pagar licencias millonarias por un núcleo ARM y por las herramientas EDA. Ahora, puedes descargar un núcleo RISC-V (como \textbf{Ibex} o \textbf{PicoRV32}) escrito en Verilog \textbf{gratis}, y usar The OpenROAD Project para implementarlo en un chip, también \textbf{gratis}.

    \item \textbf{Innovación sin Permiso:} ¿Quieres añadir una instrucción personalizada a tu procesador para acelerar una tarea específica (IA, criptografía, etc.)? Con RISC-V, puedes extender la ISA base con tus propias instrucciones y luego usar OpenROAD para construir un chip que las ejecute. Esto permite una personalización y especialización sin precedentes.

    \item \textbf{Educación e Investigación:} Universidades y estudiantes ahora pueden aprender sobre arquitectura de computadoras y diseño de chips con un flujo completo y práctico. Pueden diseñar, implementar y analizar un procesador real, algo que antes estaba reservado para empresas con grandes presupuestos.
\end{enumerate}

\subsection{¿Cómo se ve en la práctica?}

Un flujo de trabajo típico para un diseñador que usa este ecosistema sería:

\begin{enumerate}
    \item \textbf{Obtener un Núcleo:} Descargas el código Verilog de un núcleo RISC-V de código abierto desde GitHub (por ejemplo, el núcleo \textbf{Ibex} de lowRISC).
    \item \textbf{Crear un SoC:} Integras ese núcleo con otros componentes (periféricos) que también pueden ser de código abierto, como controladores de memoria, UARTs, SPI, etc. Este conjunto forma tu 'System-on-Chip'.
    \item \textbf{Preparar el Flujo:} Creas los archivos de configuración para OpenROAD (\texttt{config.tcl}, \texttt{constraints.sdc}) apuntando a los archivos Verilog de tu SoC.
    \item \textbf{Ejecutar OpenROAD:} Lanzas el flujo automatizado que sintetiza, coloca y enruta tu diseño completo de SoC basado en RISC-V.
    \item \textbf{Obtener el GDSII:} Al final, obtienes el archivo \texttt{gds} de tu propio chip personalizado, listo para ser enviado a fabricar.
\end{enumerate}

En resumen, \textbf{RISC-V proporciona el 'qué' (el diseño lógico del procesador) y OpenROAD proporciona el 'cómo' (las herramientas para transformar ese diseño en un layout físico)}. Juntos, están cambiando fundamentalmente quién puede diseñar chips y qué tipo de chips se pueden crear.

\newpage
\section{El Objetivo Final: System-on-Chip (SoC)}

\begin{pregunta}
¿Qué es un System-on-Chip (SoC) y cuáles son sus componentes típicos además del procesador?
\end{pregunta}

Entender qué es un System-on-Chip (SoC) es crucial, porque es el objetivo final de la mayoría de los flujos de diseño modernos, incluyendo el que utiliza RISC-V y OpenROAD. Un SoC no es solo un procesador; es, como su nombre indica, un \textbf{sistema completo integrado en un único chip}.

\subsection{¿Qué es un SoC?}

Un \textbf{System-on-Chip (SoC)} es un circuito integrado que reúne todos o la mayoría de los componentes de un sistema electrónico completo en una sola pieza de silicio. En lugar de construir un sistema con una placa de circuito impreso (PCB) que contiene un chip para el procesador, otro para la memoria, etc., un SoC integra todas estas funciones en un único chip.

\textbf{Analogía:} Piensa en la diferencia entre un ordenador de sobremesa de los años 90 y un smartphone moderno.
\begin{itemize}
    \item El \textbf{ordenador de sobremesa} es como un sistema basado en PCB: tiene una placa base con un zócalo para la CPU, ranuras para la RAM, una tarjeta gráfica separada, una tarjeta de sonido, etc.
    \item El \textbf{smartphone} funciona gracias a un SoC: un único chip contiene el procesador principal, el procesador gráfico, la memoria, el módem para la conectividad celular, el controlador Wi-Fi/Bluetooth y mucho más.
\end{itemize}

\subsection{Componentes Típicos de un SoC}

El procesador (CPU, como un núcleo RISC-V) es el cerebro, pero necesita un equipo de especialistas a su alrededor para funcionar. Estos son los componentes más comunes que encontrarás en un SoC moderno:

\begin{enumerate}
    \item \textbf{Procesadores Especializados:}
    \begin{itemize}
        \item \textbf{GPU (Unidad de Procesamiento Gráfico):} Optimizada para renderizar gráficos.
        \item \textbf{DSP (Procesador de Señales Digitales):} Para procesar señales del mundo real (audio, video, etc.).
        \item \textbf{NPU/Acelerador de IA (Unidad de Procesamiento Neuronal):} Hardware para acelerar operaciones de redes neuronales.
    \end{itemize}

    \item \textbf{Memoria:}
    \begin{itemize}
        \item \textbf{SRAM (RAM Estática):} Memoria muy rápida integrada en el chip, usada como cachés (L1, L2).
        \item \textbf{ROM (Memoria de Solo Lectura):} Contiene el código de arranque inicial (bootloader).
        \item \textbf{Controlador de Memoria (DRAM Controller):} Una interfaz para comunicarse con la memoria RAM principal (DRAM), que suele ser un chip externo.
    \end{itemize}

    \item \textbf{Periféricos de Conectividad:}
    \begin{itemize}
        \item \textbf{Interfaces de baja velocidad:} UART, SPI, I²C (para conectar sensores y otros chips).
        \item \textbf{Interfaces de alta velocidad:} USB, Ethernet, PCIe.
        \item \textbf{Módems de Radiofrecuencia (RF):} Para Wi-Fi, Bluetooth y conectividad celular.
    \end{itemize}

    \item \textbf{Periféricos Multimedia:}
    \begin{itemize}
        \item \textbf{Controlador de Pantalla:} Genera la señal de video para una pantalla.
        \item \textbf{Codificadores/Decodificadores de Video (Codecs):} Hardware para comprimir y descomprimir video de manera eficiente.
    \end{itemize}

    \item \textbf{Buses de Interconexión (Interconnect Fabric):}
    \begin{itemize}
        \item Son las 'autopistas' dentro del chip que conectan todos los componentes. Estándares comunes son \textbf{AMBA (AXI, AHB, APB)}.
    \end{itemize}

    \item \textbf{Sistema de Reloj y Alimentación:}
    \begin{itemize}
        \item \textbf{PLLs (Phase-Locked Loops):} Generan las diferentes señales de reloj para los distintos componentes.
        \item \textbf{Unidades de Gestión de Energía:} Controlan el consumo de energía para ahorrar batería.
    \end{itemize}

    \item \textbf{Seguridad:}
    \begin{itemize}
        \item \textbf{Motores Criptográficos:} Hardware para acelerar el cifrado y descifrado.
        \item \textbf{Root of Trust (Raíz de Confianza):} Hardware para asegurar un arranque seguro del sistema.
    \end{itemize}
\end{enumerate}

\newpage
\section{Estructura de un Proyecto de SoC con RISC-V}

\begin{pregunta}
Muéstrame un ejemplo de cómo se vería la estructura de archivos de un proyecto de SoC simple basado en RISC-V.
\end{pregunta}

Ver la estructura de un proyecto real es la mejor manera de conectar todos los conceptos. Un proyecto de SoC, incluso uno simple, tiene varias partes móviles que necesitan estar bien organizadas. A continuación se muestra una estructura de directorios típica para un SoC simple que contiene un núcleo RISC-V, una RAM en chip y un periférico UART.

\subsection{Estructura de Archivos del Proyecto \texttt{mi\_soc\_riscv}}

\begin{verbatim}
mi_soc_riscv/
|
|-- doc/
|   `-- memory_map.md         # Documentacion del mapa de memoria del SoC.
|
|-- fw/ (Firmware / Software)
|   |-- hello.c               # El programa en C que se ejecutara en el procesador.
|   |-- linker.ld             # Script del enlazador para colocar el codigo en la RAM.
|   `-- Makefile              # Para compilar el C a un binario y a un .hex.
|
|-- hdl/ (Hardware Description Language - Verilog)
|   |-- soc_top.v             # Modulo principal que conecta todo el SoC.
|   |-- simple_mem.v          # Modulo de la memoria RAM en el chip.
|   |-- simple_uart.v         # Modulo del periferico UART.
|   |-- interconnect.v        # Logica para conectar la CPU a la memoria y al UART.
|   `-- picorv32/             # Submodulo Git con el codigo del nucleo PicoRV32.
|       `-- picorv32.v
|
|-- sim/ (Simulacion)
|   |-- tb_soc_top.v          # Banco de pruebas para simular el SoC completo.
|   `-- firmware.hex          # El programa compilado, que el testbench carga.
|
`-- flow/ (Flujo de Diseno Fisico - OpenROAD)
    |-- config.tcl            # Archivo de configuracion principal para OpenROAD.
    |-- constraints.sdc       # Restricciones de tiempo (frecuencia de reloj, etc.).
    `-- run_flow.sh           # Script para ejecutar el flujo completo de OpenROAD.
\end{verbatim}

\subsection{Explicación de cada Directorio}

\begin{description}
    \item[\texttt{doc/}:] Contiene la documentación. El archivo más importante es el \textbf{mapa de memoria (\texttt{memory\_map.md})}, que define en qué dirección de memoria se encuentra cada periférico (ej. RAM en \texttt{0x00000000}, UART en \texttt{0x00010000}).
    
    \item[\texttt{fw/}:] Contiene el software que se ejecutará en el chip. Se usa una cadena de herramientas (toolchain) de RISC-V para compilar el código C (\texttt{hello.c}) en un archivo hexadecimal (\texttt{firmware.hex}) que el hardware pueda entender.
    
    \item[\texttt{hdl/}:] El corazón del hardware. Contiene todo el código Verilog. El módulo \texttt{soc\_top.v} es el más importante, ya que instancia y conecta el núcleo del procesador, la memoria y todos los periféricos.
    
    \item[\texttt{sim/}:] Para verificar que el hardware funciona antes de fabricarlo. El banco de pruebas (\texttt{tb\_soc\_top.v}) instancia el SoC, le proporciona una señal de reloj, carga el \texttt{firmware.hex} en la memoria simulada y comprueba si el comportamiento es el esperado (ej. si el UART imprime el texto correcto).
    
    \item[\texttt{flow/}:] Contiene todo lo necesario para el flujo de diseño físico con OpenROAD. El archivo \texttt{config.tcl} le dice a OpenROAD qué archivos Verilog usar, qué tecnología de fabricación (PDK) y otras configuraciones. El archivo \texttt{constraints.sdc} define los objetivos de tiempo, como la frecuencia de reloj.
\end{description}

Este tipo de organización mantiene el proyecto limpio y separa claramente las diferentes responsabilidades: el diseño de hardware (\texttt{hdl}), el software que corre en él (\texttt{fw}), la verificación (\texttt{sim}) y la implementación física (\texttt{flow}).

\newpage
\section{La Interconexión del SoC: El Bus AMBA}

\begin{pregunta}
¿Qué es el bus AMBA (AXI, AHB, APB) y cómo funciona para conectar los componentes de un SoC?
\end{pregunta}

Si los componentes del SoC son los 'órganos' del sistema, el bus AMBA es el \textbf{sistema nervioso y circulatorio} que los conecta a todos. Es imposible construir un SoC sin entender cómo se comunican sus partes.

\subsection{¿Qué es AMBA?}

\textbf{AMBA (Advanced Microcontroller Bus Architecture)} es un estándar de interconexión en chip, de código abierto y libre de regalías, definido por ARM. Su propósito es estandarizar cómo se comunican los diferentes bloques (el procesador, la memoria, los periféricos) dentro de un SoC.

Aunque fue creado por ARM, se ha convertido en el \textbf{estándar de facto de la industria}. La gran mayoría de los SoCs, incluidos los basados en \textbf{RISC-V}, utilizan buses compatibles con AMBA para su interconexión interna.

AMBA no es un único bus, sino una familia de protocolos de bus, cada uno optimizado para un propósito diferente, formando una jerarquía de rendimiento.

\subsection{APB (Advanced Peripheral Bus) - La 'Calle Residencial'}
\begin{itemize}
    \item \textbf{Propósito:} Conectar periféricos de \textbf{baja velocidad y bajo consumo}.
    \item \textbf{Características:} Es el bus más sencillo de la familia, sin operaciones complejas. Su simplicidad lo hace ideal para partes del chip que deben consumir muy poca energía.
    \item \textbf{Componentes Típicos:} UART, SPI, I²C, GPIO (pines de entrada/salida de propósito general), temporizadores.
    \item \textbf{Analogía:} Es la calle residencial de tu barrio. El tráfico es lento, pero es muy simple y eficiente para llegar a las casas (los periféricos simples).
\end{itemize}

\subsection{AHB (Advanced High-performance Bus) - La 'Avenida Principal'}
\begin{itemize}
    \item \textbf{Propósito:} Para componentes que necesitan un \textbf{rendimiento medio}.
    \item \textbf{Características:} Permite operaciones más complejas como transferencias en ráfaga (bursts) y pipelining.
    \item \textbf{Componentes Típicos:} Controladores de memoria en el chip (SRAM), controladores DMA (Acceso Directo a Memoria).
    \item \textbf{Analogía:} Es la avenida principal de la ciudad. Soporta más tráfico y a mayor velocidad que una calle residencial.
\end{itemize}

\subsection{AXI (Advanced eXtensible Interface) - La 'Autopista de Alta Velocidad'}
\begin{itemize}
    \item \textbf{Propósito:} Es el bus de \textbf{más alto rendimiento}, diseñado para conectar los componentes más rápidos y exigentes del SoC.
    \item \textbf{Características:} Está basado en cinco canales independientes (dirección/datos de lectura/escritura), lo que permite un paralelismo masivo y múltiples transacciones pendientes simultáneamente.
    \item \textbf{Componentes Típicos:} El procesador (CPU), el procesador gráfico (GPU), controladores de memoria externa (DDR), interconexiones de alta velocidad como PCIe.
    \item \textbf{Analogía:} Es una autopista de múltiples carriles con carriles separados para cada dirección, permitiendo un flujo masivo de tráfico a alta velocidad.
\end{itemize}

\subsection{Arquitectura Jerárquica y Puentes (Bridges)}

Un SoC típico no usa un solo tipo de bus, sino que los combina en una jerarquía para ser eficiente. Para conectar buses de diferentes tipos, se utilizan \textbf{puentes (bridges)}. Una arquitectura común se vería así:

\begin{enumerate}
    \item El \textbf{procesador (CPU)} y la \textbf{GPU} se conectan a una \textbf{interconexión AXI} de alta velocidad. El controlador de memoria DDR también está en este bus.
    \item Para hablar con periféricos más lentos, la interconexión AXI se conecta a un \textbf{puente AXI-a-APB}.
    \item Este puente es el 'maestro' en un \textbf{bus APB}, al cual se conectan como 'esclavos' todos los periféricos de baja velocidad como el \textbf{UART}, \textbf{SPI} y \textbf{GPIO}.
\end{enumerate}

\begin{verbatim}
          +-------------+      +-----------------+
CPU Core  |             |      |                 |
(AXI) <-----> AXI Interconnect <-----> DDR Controller
          |             |      |      (AXI)      |
          +-------------+      +-----------------+
                 |
                 |
          +----------------+
          |                |
          | AXI-to-APB Bridge |
          |                |
          +----------------+
                 | (Bus APB)
  +--------------+--------------+
  |              |              |
+-------+    +-------+    +-------+
| UART  |    |  SPI  |    | GPIO  |
+-------+    +-------+    +-------+
\end{verbatim}

Esta estructura permite que el procesador realice transacciones de alta velocidad con la memoria sin ser ralentizado por el acceso lento a un periférico simple.

\newpage
\section{Ejemplo de Código Verilog: \texttt{soc\_top.v} y \texttt{interconnect.v}}

\begin{pregunta}
Muéstrame un ejemplo simple del código Verilog para el 'soc\_top.v' y el 
'interconnect.v'.
\end{pregunta}

¡Claro! Ver el código Verilog para estos dos módulos es la mejor manera de entender cómo se 'ensambla' el hardware a nivel de RTL. Este código asume la arquitectura que hemos discutido: un núcleo de CPU, una memoria RAM y un periférico UART.

\subsection{Prerrequisito: El Mapa de Memoria}

Antes de escribir el código, necesitamos un mapa de memoria. Este es el acuerdo que define qué dirección corresponde a qué componente. Para nuestro SoC simple, usaremos:
\begin{itemize}
    \item \textbf{RAM:} \texttt{0x00000000} - \texttt{0x0000FFFF} (64 KB)
    \item \textbf{UART:} \texttt{0x10000000} - \texttt{0x1000000F} (Unas pocas direcciones para sus registros)
\end{itemize}
El \texttt{interconnect} usará este mapa para dirigir el tráfico.

\subsection{1. \texttt{interconnect.v} - El Decodificador de Direcciones}

Este módulo se sitúa entre la CPU (el 'maestro') y los periféricos (los 'esclavos'). Su única función es mirar la dirección que pide la CPU y enrutar la petición al dispositivo correcto.

\begin{lstlisting}[language=Verilog, caption={interconnect.v}]
/*
 * interconnect.v
 *
 * Este modulo actua como un decodificador de direcciones simple.
 * Conecta un maestro (la CPU) a dos esclavos (RAM y UART)
 * basandose en el mapa de memoria.
 */
module interconnect (
    // --- Interfaz con el Maestro (CPU) ---
    input               cpu_mem_valid,
    output reg          cpu_mem_ready,
    input      [31:0]   cpu_mem_addr,
    input      [31:0]   cpu_mem_wdata,
    input      [3:0]    cpu_mem_wstrb,
    output reg [31:0]   cpu_mem_rdata,

    // --- Interfaz con el Esclavo 1 (RAM) ---
    output reg          ram_mem_valid,
    input               ram_mem_ready,
    output reg [31:0]   ram_mem_addr,
    output reg [31:0]   ram_mem_wdata,
    output reg [3:0]    ram_mem_wstrb,
    input      [31:0]   ram_mem_rdata,

    // --- Interfaz con el Esclavo 2 (UART) ---
    output reg          uart_mem_valid,
    input               uart_mem_ready,
    output reg [31:0]   uart_mem_addr,
    output reg [31:0]   uart_mem_wdata,
    output reg [3:0]    uart_mem_wstrb,
    input      [31:0]   uart_mem_rdata
);

    // Logica de decodificacion de direccion
    localparam ADDR_MASK_RAM  = 32'h0000_0000;
    localparam ADDR_MASK_UART = 32'h1000_0000;

    wire select_ram  = (cpu_mem_addr[31:16] == ADDR_MASK_RAM[31:16]);
    wire select_uart = (cpu_mem_addr[31:28] == ADDR_MASK_UART[31:28]);

    always @(*) begin
        // Por defecto, ningun esclavo esta activo
        ram_mem_valid  = 1'b0;
        uart_mem_valid = 1'b0;

        // Pasa las senales de la CPU al esclavo seleccionado
        if (cpu_mem_valid) begin
            if (select_ram) begin
                ram_mem_valid = 1'b1;
                ram_mem_addr  = cpu_mem_addr;
                ram_mem_wdata = cpu_mem_wdata;
                ram_mem_wstrb = cpu_mem_wstrb;
            end else if (select_uart) begin
                uart_mem_valid = 1'b1;
                uart_mem_addr  = cpu_mem_addr;
                uart_mem_wdata = cpu_mem_wdata;
                uart_mem_wstrb = cpu_mem_wstrb;
            end
        end
    end

    // Logica para enrutar las senales de vuelta a la CPU
    always @(*) begin
        cpu_mem_ready = (select_ram && ram_mem_ready) || (select_uart && uart_mem_ready);

        if (select_ram) begin
            cpu_mem_rdata = ram_mem_rdata;
        end else if (select_uart) begin
            cpu_mem_rdata = uart_mem_rdata;
        end else begin
            cpu_mem_rdata = 32'hDEADBEEF; // Valor de error
        end
    end

endmodule
\end{lstlisting}

\subsection{2. \texttt{soc\_top.v} - El Ensamblaje Final}

Este es el módulo de más alto nivel. No contiene mucha lógica, su principal trabajo es \textbf{instanciar} todos los demás módulos (CPU, RAM, UART, Interconnect) y \textbf{conectar sus puertos} entre sí.

\begin{lstlisting}[language=Verilog, caption={soc\_top.v}]
/*
 * soc_top.v
 *
 * Modulo de mas alto nivel para nuestro SoC simple.
 * Instancia y conecta la CPU, la memoria, el UART y el interconnect.
 */
module soc_top (
    input  clk,
    input  rst,
    output uart_tx,
    input  uart_rx
);

    // Wires para las interfaces de memoria
    wire cpu_mem_valid, cpu_mem_ready;
    wire [31:0] cpu_mem_addr, cpu_mem_wdata, cpu_mem_rdata;
    wire [3:0]  cpu_mem_wstrb;
    // ... (declaracion de wires para ram y uart)

    // 1. Instancia del Nucleo de la CPU (PicoRV32)
    picorv32 cpu (
        .clk(clk), .resetn(~rst),
        .mem_valid(cpu_mem_valid), .mem_ready(cpu_mem_ready),
        .mem_addr(cpu_mem_addr), .mem_wdata(cpu_mem_wdata),
        .mem_wstrb(cpu_mem_wstrb), .mem_rdata(cpu_mem_rdata)
    );

    // 2. Instancia del Interconnect
    interconnect bus_logic (
        .cpu_mem_valid(cpu_mem_valid), .cpu_mem_ready(cpu_mem_ready),
        .cpu_mem_addr(cpu_mem_addr), .cpu_mem_wdata(cpu_mem_wdata),
        .cpu_mem_wstrb(cpu_mem_wstrb), .cpu_mem_rdata(cpu_mem_rdata),

        .ram_mem_valid(ram_mem_valid), .ram_mem_ready(ram_mem_ready),
        /* ... (conexiones a la RAM) ... */

        .uart_mem_valid(uart_mem_valid), .uart_mem_ready(uart_mem_ready)
        /* ... (conexiones al UART) ... */
    );

    // 3. Instancia de la Memoria RAM
    simple_mem ram (
        .clk(clk), .mem_valid(ram_mem_valid), .mem_ready(ram_mem_ready),
        .mem_addr(ram_mem_addr[15:2]), .mem_wdata(ram_mem_wdata),
        .mem_wstrb(ram_mem_wstrb), .mem_rdata(ram_mem_rdata)
    );

    // 4. Instancia del Periferico UART
    simple_uart uart (
        .clk(clk), .rst(rst),
        .mem_valid(uart_mem_valid), .mem_ready(uart_mem_ready),
        .mem_addr(uart_mem_addr[3:2]), .mem_wdata(uart_mem_wdata),
        .mem_wstrb(uart_mem_wstrb), .mem_rdata(uart_mem_rdata),
        .tx(uart_tx), .rx(uart_rx)
    );

endmodule
\end{lstlisting}

\newpage
\section{Optimizando el Flujo de Datos: El Controlador DMA}

\begin{pregunta}
¿Qué es un controlador DMA y qué papel juega en un SoC?
\end{pregunta}

¡Excelente pregunta! El controlador DMA es uno de los componentes más importantes y a la vez menos 'visibles' de un SoC. Es un especialista en eficiencia que permite que todo el sistema funcione mucho más rápido.

Si la CPU es el 'cerebro' pensante del SoC, el \textbf{controlador DMA es el 'sistema logístico' o el 'equipo de mudanzas'}.

\subsection{¿Qué es un Controlador DMA?}

\textbf{DMA} son las siglas de \textbf{Direct Memory Access} (Acceso Directo a Memoria). Un controlador DMA es un bloque de hardware especializado cuya única función es \textbf{mover grandes cantidades de datos de un lugar a otro dentro del SoC, sin la intervención de la CPU}.

\subsection{El Problema: La CPU como Cuello de Botella}

Imagina que un periférico de red acaba de recibir un paquete de datos de 1500 bytes. Sin un DMA, la CPU tendría que leer cada byte del periférico y escribirlo en la memoria, uno por uno, 1500 veces. Durante todo este tiempo, la CPU, un recurso increíblemente potente, está ocupada haciendo una tarea repetitiva y 'tonta' de copiar y pegar. Es un enorme desperdicio de ciclos de computación.

\subsection{La Solución: Delegar en el Especialista (DMA)}

Con un controlador DMA, el proceso es mucho más eficiente:
\begin{enumerate}
    \item La CPU configura el controlador DMA. Le da una orden muy simple:
    \begin{itemize}
        \item \textbf{Dirección de Origen:} El registro del periférico donde están los datos.
        \item \textbf{Dirección de Destino:} La dirección en la memoria RAM donde se deben guardar.
        \item \textbf{Tamaño:} 1500 bytes.
    \end{itemize}
    \item La CPU le dice al DMA: '¡Empieza!'. Y \textbf{la CPU queda libre inmediatamente} para hacer otras tareas importantes.
    \item El controlador DMA toma el control del bus del sistema (se convierte en un 'bus master') y realiza la transferencia completa directamente.
    \item Una vez que la transferencia ha terminado, el controlador DMA envía una única señal a la CPU (una \textbf{interrupción}) para decirle: 'He terminado'.
\end{enumerate}

\subsection{El Papel que Juega en un SoC}

El controlador DMA es un componente fundamental para el rendimiento en cualquier SoC que maneje grandes volúmenes de datos.

\begin{enumerate}
    \item \textbf{Libera a la CPU:} Su función principal es descargar a la CPU de las tareas de transferencia de datos, permitiendo que el procesador se dedique a la computación, lo que aumenta drásticamente el rendimiento general del sistema.

    \item \textbf{Permite el Alto Rendimiento de Periféricos:} Es esencial para periféricos que generan o consumen datos a alta velocidad. Algunos ejemplos son:
    \begin{itemize}
        \item \textbf{Redes:} Controladores Ethernet y Wi-Fi.
        \item \textbf{Almacenamiento:} Controladores de SSDs (NVMe/PCIe).
        \item \textbf{Gráficos:} Mover texturas y datos de vértices a la memoria de la GPU.
        \item \textbf{Audio/Video:} Transferir flujos de audio o fotogramas de video desde un códec o una cámara.
        \item \textbf{Adquisición de Datos:} Recoger muestras de un Convertidor Analógico-Digital (ADC) a alta velocidad.
    \end{itemize}

    \item \textbf{Posición en la Arquitectura del SoC:} Como el DMA necesita mover datos rápidamente, debe ser un 'ciudadano de primera clase' en el bus del sistema. Por lo general, se conecta a un bus de alto rendimiento como \textbf{AXI} o \textbf{AHB}, donde puede actuar como un 'maestro' del bus, al igual que la CPU.
\end{enumerate}

\subsection{Analogía Final}

Piensa en la CPU como el chef principal de una cocina muy ocupada.
\begin{itemize}
    \item \textbf{Sin DMA:} El chef tiene que ir personalmente al almacén a por cada ingrediente (cada byte), uno por uno, mientras las sartenes se queman.
    \item \textbf{Con DMA:} El chef le da una lista de ingredientes 
    (la configuración del DMA) a un ayudante de cocina (el controlador DMA) y 
    le dice 'tráeme todo esto a mi mesa'. Mientras el ayudante va al almacén y 
    vuelve, el chef puede seguir cocinando y preparando los platos. La cocina es 
    mucho más productiva.
\end{itemize}

\newpage
\section{Verificación Funcional: El Banco de Pruebas (Testbench)}

\begin{pregunta}
¿Cómo se simula este SoC con un testbench para verificar que el firmware 'Hola, Mundo' funciona?
\end{pregunta}

Este es el paso final y uno de los más importantes del ciclo de diseño: la \textbf{verificación funcional}. Antes de fabricar un chip, debemos estar lo más seguros posible de que el hardware hace lo que se supone que debe hacer. Para ello, creamos un \textbf{banco de pruebas} o \textbf{testbench}.

\subsection{¿Qué es un Testbench?}

Un testbench es un módulo de Verilog que \textbf{no es sintetizable}, es decir, no se convierte en hardware. Su único propósito es existir en el entorno de simulación para:
\begin{enumerate}
    \item Instanciar nuestro diseño (el 'Device Under Test' o DUT), que en este caso es \texttt{soc\_top}.
    \item Proporcionarle las señales de entrada necesarias, como el reloj (\texttt{clk}) y el reset (\texttt{rst}).
    \item Cargar el programa (\texttt{firmware.hex}) en la memoria del SoC.
    \item Monitorizar las salidas (como el pin \texttt{uart\_tx}) y comprobar si son las esperadas.
    \item Informar si la prueba fue un éxito o un fracaso.
\end{enumerate}

\subsection{\texttt{tb\_soc\_top.v} - El Banco de Pruebas para Nuestro SoC}

Aquí tienes un ejemplo de cómo se vería el archivo \texttt{sim/tb\_soc\_top.v}.

\begin{lstlisting}[language=Verilog, caption={tb\_soc\_top.v}]
/*
 * tb_soc_top.v
 *
 * Banco de pruebas para nuestro SoC simple.
 * 1. Genera el reloj y el reset.
 * 2. Carga el firmware 'firmware.hex' en la memoria RAM.
 * 3. Monitoriza la salida del UART y la imprime en la consola.
 */
`timescale 1ns / 1ps

module tb_soc_top;

    localparam CLK_PERIOD = 10; // 10 ns -> 100 MHz

    reg clk;
    reg rst;
    wire uart_tx;

    // 1. Instancia del DUT (Device Under Test)
    soc_top dut (
        .clk(clk),
        .rst(rst),
        .uart_tx(uart_tx),
        .uart_rx(1'b1) // No se usa, se conecta a 1
    );

    // 2. Generacion del Reloj
    initial begin
        clk = 0;
        forever #(CLK_PERIOD / 2) clk = ~clk;
    end

    // 3. Generacion del Reset y Carga del Firmware
    initial begin
        $display('Inicio de la simulacion.');
        rst = 1;
        #(CLK_PERIOD * 10);
        rst = 0;
        
        // Cargar el firmware en la memoria del SoC
        // La ruta apunta a la instancia de la RAM dentro del DUT.
        $display('Cargando firmware.hex en la memoria...');
        $readmemh('firmware.hex', dut.ram.memory);
        
        $display('Firmware cargado. La CPU comenzara a ejecutar.');

        // Terminar la simulacion despues de un tiempo
        #500000;
        $display('Timeout de simulacion alcanzado. Terminando.');
        $finish;
    end

    // 4. Monitor y Decodificador del UART (simplificado)
    // Una tarea real seria mas compleja, pero esto muestra la idea.
    task automatic display_uart_char;
        // ... (logica para decodificar 8 bits desde uart_tx)
        // $write('%c', received_char);
    endtask

endmodule
\end{lstlisting}

\newpage
\section{La Comunicación Asíncrona: Interrupciones}

\begin{pregunta}
¿Qué son las interrupciones y cómo las maneja un procesador como RISC-V?
\end{pregunta}

Hemos hablado de cómo la CPU se comunica con los periféricos (a través de buses) y cómo el DMA puede mover datos por sí solo. Pero, ¿cómo le avisa un periférico a la CPU que ha ocurrido un evento importante \textit{ahora mismo}? (por ejemplo, '¡He recibido un nuevo byte por el UART!' o '¡El temporizador ha llegado a cero!').

La respuesta son las \textbf{interrupciones}. Son el mecanismo que permite al hardware externo demandar la atención de la CPU de forma asíncrona.

\subsection{¿Qué es una Interrupción?}

Una interrupción es una señal que envía un periférico a la CPU para indicar que ha ocurrido un evento que requiere atención inmediata. Cuando la CPU recibe esta señal, \textbf{detiene (interrumpe) su tarea actual}, atiende la petición del periférico y, una vez terminada, \textbf{reanuda la tarea original} exactamente donde la dejó.

\textbf{Analogía:} Imagina que estás leyendo un libro (la tarea principal de la CPU) y de repente suena el timbre de la puerta (la interrupción). Pones un marcapáginas en la página que estás leyendo (guardas el contexto), vas a atender la puerta (ejecutas la rutina de servicio de interrupción), y cuando terminas, vuelves al libro, buscas el marcapáginas y sigues leyendo como si nada hubiera pasado.

\subsection{El Flujo de una Interrupción}

El proceso sigue una secuencia bien definida:
\begin{enumerate}
    \item \textbf{Petición de Interrupción:} Un periférico (ej. un UART) activa su línea de interrupción para señalar un evento (ej. ha llegado un dato).
    \item \textbf{Detección por el Controlador:} La señal llega a un \textbf{Controlador de Interrupciones}, un hardware especializado que gestiona las peticiones de múltiples periféricos. Este controlador prioriza las interrupciones y notifica a la CPU.
    \item \textbf{Pausa y Guardado de Contexto:} La CPU termina la instrucción que está ejecutando, y luego guarda su estado actual. Como mínimo, guarda el \textbf{Contador de Programa (PC)} en un registro especial. A menudo también guarda otros registros en una zona de memoria llamada la \textbf{pila (stack)}.
    \item \textbf{Salto a la Rutina de Servicio (ISR):} La CPU salta a una dirección de memoria predefinida que está asociada con esa interrupción específica. El código que se encuentra en esa dirección se llama \textbf{Rutina de Servicio de Interrupción (ISR - Interrupt Service Routine)}.
    \item \textbf{Ejecución de la ISR:} La ISR se ejecuta. Su código está diseñado para 'atender' al periférico que causó la interrupción (por ejemplo, leer el byte del registro del UART y guardarlo en un buffer).
    \item \textbf{Retorno y Restauración de Contexto:} Una vez que la ISR termina, ejecuta una instrucción especial de 'retorno de interrupción' (en RISC-V, es \texttt{mret}). Esta instrucción le dice a la CPU que recupere el estado que guardó en el paso 3 (especialmente el PC) y continúe con el programa original.
\end{enumerate}

Para el programa principal, esta interrupción fue completamente transparente (aparte de una pequeña pérdida de tiempo).

\subsection{Interrupciones en el Ecosistema RISC-V}

La especificación RISC-V define cómo el \textit{núcleo} de la CPU debe reaccionar a una interrupción, pero no define cómo se gestionan las interrupciones de \textit{múltiples} periféricos externos. Para estandarizar esto, el ecosistema RISC-V utiliza un componente llamado \textbf{PLIC (Platform-Level Interrupt Controller)}.

\begin{description}
    \item[PLIC (Platform-Level Interrupt Controller):] Es un bloque de hardware que se sitúa entre los periféricos del SoC y el núcleo de la CPU. Sus funciones son:
    \begin{itemize}
        \item \textbf{Recibir} las peticiones de interrupción de todos los periféricos (UART, SPI, DMA, etc.).
        \item \textbf{Priorizar} las interrupciones. Si dos periféricos piden atención al mismo tiempo, el PLIC decide cuál es más importante basándose en su configuración.
        \item \textbf{Enmascarar} interrupciones. Permite al software habilitar o deshabilitar interrupciones de periféricos específicos.
        \item \textbf{Notificar} a la CPU a través de una única línea de interrupción externa.
        \item \textbf{Informar} a la CPU sobre \textit{cuál} periférico necesita atención, para que la CPU sepa qué ISR debe ejecutar.
    \end{itemize}
\end{description}

En resumen, cuando construyes un SoC basado en RISC-V, además del núcleo de la CPU y los periféricos, casi siempre incluirás una instancia de un PLIC para gestionar las interrupciones de manera ordenada y eficiente.

\newpage
\section{El Puente a la Realidad Física: El PDK}

\begin{pregunta}
Explica qué es un PDK (Process Design Kit) y por qué es crucial para el flujo de OpenROAD.
\end{pregunta}

Hasta ahora, hemos hablado de diseñar en Verilog y usar OpenROAD para generar un layout. Pero, ¿cómo sabe OpenROAD las reglas físicas de la fábrica (foundry) a la que enviaremos el chip? ¿Cómo sabe qué tan gruesos deben ser los cables, qué tan separadas deben estar las compuertas o qué capas de metal están disponibles?

La respuesta es el \textbf{PDK (Process Design Kit)}.

\subsection{¿Qué es un PDK?}

Un PDK es el 'manual de reglas y catálogo de piezas' que proporciona una fundición de semiconductores específica (como SkyWater, GlobalFoundries, o TSMC). Es una colección de archivos de datos que describe detalladamente un proceso de fabricación de semiconductores.

\textbf{Analogía:} Si OpenROAD es el equipo de arquitectos e ingenieros que diseñan un rascacielos, el PDK es el conjunto de códigos de construcción locales, el catálogo de vigas de acero disponibles, y las especificaciones del hormigón que deben usar. Sin el PDK, los arquitectos no sabrían qué materiales usar ni qué reglas seguir, y el edificio se derrumbaría.

OpenROAD es una herramienta de propósito general; no sabe nada sobre un proceso de fabricación específico. El PDK le proporciona todo el conocimiento tecnológico necesario para transformar un diseño lógico genérico en un layout físico que se pueda fabricar con éxito.

\subsection{Componentes Clave de un PDK}

Un PDK es un conjunto de librerías y archivos, donde cada uno tiene un rol específico en el flujo de diseño:

\begin{enumerate}
    \item \textbf{Librerías de Celdas Estándar (Standard Cells):} Son las piezas de LEGO básicas del diseño digital. El PDK proporciona múltiples 'vistas' de cada celda (AND, OR, NOT, Flip-Flop, etc.):
    \begin{itemize}
        \item \textbf{Vista Lógica (\texttt{.lib}):} Un archivo de texto que describe la función lógica y, lo más importante, las \textbf{características de tiempo} de cada celda (retrasos de propagación, tiempos de subida/bajada). Es crucial para la síntesis y el análisis de tiempo estático (STA).
        \item \textbf{Vista Física Abstracta (\texttt{.lef}):} Describe la forma física abstracta de la celda: su tamaño, la ubicación de sus pines y las zonas donde no se puede enrutar. Es esencial para el \textit{placement} y el \textit{routing}.
        \item \textbf{Vista de Layout Completa (\texttt{.gds}):} El diseño geométrico real de los transistores y conexiones dentro de la celda. Se usa para construir el GDSII final.
    \end{itemize}

    \item \textbf{Reglas de Diseño (DRC - Design Rule Checking):} Es el 'código de construcción' del proceso. Define cientos de reglas geométricas, como el ancho mínimo de un cable de metal, la separación mínima entre dos cables, el tamaño de las vías, etc. El \textit{router} de OpenROAD debe obedecer estas reglas para que el chip sea fabricable.

    \item \textbf{Archivo de Tecnología (\texttt{.tech.lef}):} Define las propiedades de las capas de interconexión. Especifica cuántas capas de metal hay disponibles, su grosor, su resistencia y capacitancia por unidad de longitud, etc. Esta información es vital para el \textit{router} y para calcular los retrasos del cableado.

    \item \textbf{Reglas de LVS (Layout vs. Schematic):} Reglas para una herramienta de verificación que comprueba si el layout geométrico final corresponde exactamente al netlist lógico original, asegurando que no se hayan introducido errores durante el diseño físico.
\end{enumerate}

\subsection{El Rol del PDK en OpenROAD}

El PDK es el \textbf{parámetro de entrada fundamental} para todo el flujo de OpenROAD. Cuando ejecutas el flujo, una de las primeras cosas que especificas es qué PDK usar (por ejemplo, \texttt{sky130}). A partir de ahí:
\begin{itemize}
    \item La \textbf{Síntesis} usa los archivos \texttt{.lib} para elegir las celdas que mejor cumplan los objetivos de tiempo.
    \item El \textbf{Floorplanning} y el \textbf{Placement} usan los archivos \texttt{.lef} para saber el tamaño de las celdas y cómo colocarlas.
    \item El \textbf{Routing} usa el \texttt{.tech.lef} para saber qué capas de metal puede usar y las reglas de DRC para asegurarse de que el cableado es válido.
    \item El \textbf{Análisis de Tiempo Estático (STA)} usa los \texttt{.lib} y la información de resistencia/capacitancia del \texttt{.tech.lef} para calcular los retrasos totales y verificar si el chip funcionará a la frecuencia deseada.
\end{itemize}

El auge de los PDKs de código abierto, como el \textbf{SkyWater 130nm (sky130)} patrocinado por Google, es lo que ha hecho posible tener un flujo de diseño de chips \textbf{completamente de código abierto}, desde el Verilog hasta el GDSII, democratizando el acceso a la fabricación de silicio.

\newpage
\section{El Gran Desafío: Timing Closure}

\begin{pregunta}
¿Qué es el 'timing closure' y por qué es uno de los mayores desafíos en el diseño físico de chips?
\end{pregunta}

Si el PDK es el 'manual de reglas', el 'timing closure' es el examen final que determina si tu diseño se gradúa para la fabricación. Es, sin duda, uno de los mayores dolores de cabeza para los ingenieros de diseño físico.

\subsection{¿Qué es el 'Timing Closure'?}

El \textbf{'Timing Closure'} (cierre de temporización) es el proceso iterativo de modificar un diseño de chip hasta que \textbf{cumple con todas sus restricciones de tiempo}. En pocas palabras, es el acto de asegurarse de que todas las señales dentro del chip llegan a donde tienen que ir en el momento correcto, ni demasiado tarde ni demasiado pronto.

El objetivo principal es que el chip funcione de manera fiable a la frecuencia de reloj deseada. Para ello, se deben satisfacer dos condiciones fundamentales en cada ruta entre dos registros (flip-flops):

\begin{enumerate}
    \item \textbf{Setup Time (Tiempo de Preparación):} La señal debe llegar al segundo registro y estar estable un cierto tiempo \textit{antes} de que llegue el siguiente pulso de reloj. Una \textbf{violación de setup} ocurre si la señal llega \textbf{demasiado tarde}.
    \item \textbf{Hold Time (Tiempo de Mantenimiento):} La señal debe permanecer estable un cierto tiempo \textit{después} de que ha llegado el pulso de reloj. Una \textbf{violación de hold} ocurre si la señal de la siguiente operación llega \textbf{demasiado rápido}.
\end{enumerate}

'Alcanzar el timing closure' significa que el diseño tiene \textbf{cero violaciones de setup y hold} en todas las millones de rutas posibles dentro del chip.

\subsection{¿Por Qué es uno de los Mayores Desafíos?}

Alcanzar el timing closure es extremadamente difícil porque cada paso del diseño físico afecta al tiempo, y a menudo, arreglar un problema crea otro en otro lugar.

\begin{enumerate}
    \item \textbf{Interdependencia de las Etapas (El Bucle Infernal):} La ubicación de las celdas (\textit{placement}) afecta la longitud de los cables, la longitud de los cables (\textit{routing}) afecta el retraso, y el retraso total (\textit{timing analysis}) te obliga a cambiar la ubicación de nuevo. Este ciclo de \texttt{Place -> Route -> Analyze -> Repeat} es el corazón del desafío.

    \item \textbf{Soluciones Contradictorias:}
    \begin{itemize}
        \item Para arreglar una \textbf{violación de setup} (ruta lenta), puedes usar celdas más grandes y potentes. Pero estas ocupan más espacio y pueden causar congestión.
        \item Para arreglar una \textbf{violación de hold} (ruta rápida), tienes que añadir retraso insertando celdas 'buffer'. Pero estos buffers también ocupan espacio y consumen energía.
    \end{itemize}
    Arreglar un tipo de violación puede fácilmente causar el otro tipo en una ruta diferente.

    \item \textbf{El Problema de los 'Corners' (PVT):} Un chip debe funcionar en un rango de \textbf{P}roceso, \textbf{V}oltaje y \textbf{T}emperatura. El timing closure debe lograrse en todos los escenarios extremos ('corners'):
    \begin{itemize}
        \item \textbf{Worst-Case (lento):} Bajo voltaje, alta temperatura. Aquí es donde ocurren las violaciones de \textbf{setup}.
        \item \textbf{Best-Case (rápido):} Alto voltaje, baja temperatura. Aquí es donde ocurren las violaciones de \textbf{hold}.
    \end{itemize}
    Optimizar para un 'corner' puede empeorar el timing en otro.

    \item \textbf{Escala y Complejidad:} En los chips modernos con miles de millones de transistores, hay cientos de millones de rutas de tiempo que deben ser analizadas y optimizadas simultáneamente, lo que representa un desafío computacional masivo.
\end{enumerate}

En resumen, el timing closure es un problema de optimización multidimensional y altamente restringido. Las herramientas EDA como OpenROAD deben equilibrar de forma inteligente el tiempo, el área, la energía y la fabricabilidad para encontrar una solución que funcione.

\newpage
\section{Hoja de Ruta del Proyecto: De un SoC a un Sistema Paralelo}

\begin{pregunta}
Ok Gemini, por ahora solo quiero diseñar, simular y dejar listo para fabricación un IC basado en RISC-V. Luego me gustaría modificar el IC para que pueda trabajar en paralelo con otros IC's igualmente basados en RISC-V.
\end{pregunta}

¡Excelente! Este es un plan de proyecto práctico y ambicioso. A continuación, se detalla una hoja de ruta en dos fases para lograr tus objetivos.

\subsection{Fase 1: Diseño y Fabricación de un SoC RISC-V Autónomo}

El objetivo aquí es crear un chip funcional y autocontenido, el "Hola, Mundo" del silicio.

\begin{enumerate}
    \item \textbf{Definición y Arquitectura:} Diseñar un microcontrolador mínimo.
    \begin{itemize}
        \item \textbf{CPU:} Un núcleo RISC-V simple como \textbf{PicoRV32}.
        \item \textbf{Memoria:} Una pequeña RAM en bloque para el firmware.
        \item \textbf{Periférico:} Un \textbf{UART} para enviar texto a un ordenador.
        \item \textbf{Entregable:} Un mapa de memoria documentado.
    \end{itemize}
    
    \item \textbf{Desarrollo del Hardware (RTL):} Describir la arquitectura en Verilog.
    \begin{itemize}
        \item Crear el módulo \texttt{soc\_top.v} que instancia y conecta la CPU, la RAM y el UART a través de un bus simple.
        \item \textbf{Entregable:} El directorio \texttt{hdl/} con todo el código Verilog.
    \end{itemize}

    \item \textbf{Desarrollo del Firmware:} Escribir el programa que se ejecutará en el chip.
    \begin{itemize}
        \item Escribir un \texttt{hello.c} que imprima "¡Hola, Mundo!" por el UART.
        \item Usar una toolchain de RISC-V para compilarlo a un archivo \texttt{firmware.hex}.
        \item \textbf{Entregable:} El archivo \texttt{firmware.hex}.
    \end{itemize}

    \item \textbf{Verificación Funcional (Simulación):} Asegurarse de que todo funciona.
    \begin{itemize}
        \item Escribir un testbench en Verilog (\texttt{tb\_soc\_top.v}) que cargue el \texttt{firmware.hex} en la RAM simulada.
        \item Ejecutar la simulación y verificar que se imprime el mensaje correcto.
        \item \textbf{Entregable:} Una simulación exitosa.
    \end{itemize}

    \item \textbf{Diseño Físico (RTL-a-GDSII):} Usar OpenROAD para crear el layout.
    \begin{itemize}
        \item Configurar y ejecutar el flujo de OpenROAD apuntando al PDK deseado (ej. \texttt{sky130}).
        \item Iterar hasta lograr el \textit{timing closure} y no tener errores de DRC.
        \item \textbf{Entregable:} El archivo \texttt{.gds} final, listo para la fábrica.
    \end{itemize}
\end{enumerate}

\subsection{Fase 2: Modificación para Trabajo en Paralelo}

El objetivo es añadir la capacidad de comunicación inter-chip.

\begin{enumerate}
    \item \textbf{Definir el Protocolo de Comunicación:} La opción más sencilla y recomendada es usar un bus \textbf{SPI (Serial Peripheral Interface)}, con una arquitectura Maestro-Esclavo.
    
    \item \textbf{Modificar el Hardware (RTL):}
    \begin{itemize}
        \item Añadir un módulo de periférico SPI al diseño del SoC.
        \item Conectarlo al bus interno y actualizar el mapa de memoria.
        \item Exponer los pines del SPI (\texttt{sclk, mosi, miso, cs}) en la interfaz del chip.
    \end{itemize}

    \item \textbf{Modificar el Firmware:}
    \begin{itemize}
        \item Escribir drivers en C para controlar el periférico SPI.
        \item Crear dos firmwares: uno para el chip \textbf{Maestro} (que envía comandos) y otro para los chips \textbf{Esclavos} (que reciben comandos, actúan y responden).
    \end{itemize}

    \item \textbf{Verificación y Diseño Físico:}
    \begin{itemize}
        \item Crear un nuevo testbench que instancie \textit{múltiples} SoCs (un maestro, uno o más esclavos) y conecte sus pines SPI para verificar la comunicación.
        \item Volver a ejecutar el flujo de OpenROAD con el diseño actualizado para generar el nuevo archivo \texttt{.gds}.
    \end{itemize}
\end{enumerate}

\newpage
\section{Estructura de una Propuesta de Financiamiento}

\begin{pregunta}
Excelente Gemini! Elaboraré un proyecto para presentarlo a las autoridades de mi Universidad para solicitar financiamiento. Espero contar contigo siempre para llevar a cabo este proyecto.
\end{pregunta}

¡Por supuesto! Es una iniciativa fantástica. Para empezar con buen pie, una propuesta de financiamiento sólida es clave. A continuación se presenta una estructura sugerida para dicho documento, basada en todo lo que hemos discutido.

\subsection{Estructura Sugerida para la Propuesta}

\begin{enumerate}
    \item \textbf{Título del Proyecto:}
    \begin{itemize}
        \item \textit{Ej: "Desarrollo de un Sistema Multi-núcleo Paralelo basado en la Arquitectura Abierta RISC-V y Herramientas EDA Open-Source".}
    \end{itemize}

    \item \textbf{Resumen Ejecutivo:}
    \begin{itemize}
        \item Un párrafo conciso que describa el proyecto, su objetivo (diseñar, verificar y preparar para fabricar un SoC, y luego extenderlo para computación paralela) y su importancia para la universidad.
    \end{itemize}

    \item \textbf{Justificación e Impacto (El "Porqué"):}
    \begin{itemize}
        \item \textbf{Innovación Académica:} Posicionar a la universidad a la vanguardia de la revolución del hardware de código abierto.
        \item \textbf{Formación de Talento:} Capacitar a los estudiantes en habilidades de diseño de VLSI altamente demandadas, utilizando herramientas modernas y accesibles.
        \item \textbf{Reducción de Costos:} Demostrar cómo el uso de herramientas como OpenROAD y estándares como RISC-V elimina las barreras de licencias millonarias.
        \item \textbf{Potencial de Investigación:} Abrir la puerta a futuras investigaciones en arquitecturas de computadoras personalizadas, aceleradores de IA, seguridad de hardware, etc.
    \end{itemize}

    \item \textbf{Objetivos del Proyecto:}
    \begin{itemize}
        \item Detallar la "Hoja de Ruta" definida previamente, con los objetivos claros para la Fase 1 (SoC autónomo) y la Fase 2 (sistema paralelo).
    \end{itemize}

    \item \textbf{Metodología:}
    \begin{itemize}
        \item Describir los pasos técnicos: Diseño RTL en Verilog, desarrollo de firmware en C, verificación con Icarus Verilog, y diseño físico con The OpenROAD Project usando el PDK \texttt{sky130}.
    \end{itemize}

    \item \textbf{Recursos Necesarios:}
    \begin{itemize}
        \item \textbf{Software:} Destacar que es 100\% de código abierto.
        \item \textbf{Hardware:} Servidores con suficiente RAM y CPU para el flujo de OpenROAD.
        \item \textbf{Costos de Fabricación:} Mencionar la viabilidad económica a través de programas de "shuttle" de bajo costo como \textbf{Tiny Tapeout} o los de Google/SkyWater.
    \end{itemize}

    \item \textbf{Cronograma y Entregables:}
    \begin{itemize}
        \item Un plan de tiempo estimado para cada fase y los entregables concretos (código fuente, informes de simulación, archivo GDSII, etc.).
    \end{itemize}

    \item \textbf{Resultados Esperados:}
    \begin{itemize}
        \item El diseño de un chip fabricable, material educativo para futuros cursos, y potencialmente una publicación académica.
    \end{itemize}
\end{enumerate}

\end{document}
